%!TEX program=xelatex
% Сборка из каталога Course_work:
%   xelatex coursework.tex && bibtex coursework && xelatex coursework.tex && xelatex coursework.tex
% На титульном листе ожидается файл img/msu.png (положите логотип ВМК / МГУ).

\documentclass[12pt]{article}
\usepackage{example}

\usepackage{hyperref}
\usepackage{amsmath}
\usepackage{caption}
\usepackage{subcaption}
\usepackage{graphicx}
\usepackage{booktabs}

\usepackage[utf8]{inputenc}

%% --- Титульный лист ---
\affiliation{Кафедра математической физики}
\author{Певцов Артём Алексеевич}
\title{Восстановление метрической карты глубины по монокулярному предсказанию и разрежённому приору: методы калибровки и экспериментальное исследование}
\worktype{КУРСОВАЯ РАБОТА}
\supervisor{уч. степень, звание\\И.О.\,Фамилия}
\date{Москва, \the\year{}}

\newfontfamily{\cyrillicfonttt}{Courier New}

\begin{document}

\maketitle
\maketoc

\section{Введение}

Оценка глубины по монокулярному изображению относится к числу фундаментальных задач машинного зрения; её решения используются при навигации роботов, построении трёхмерных моделей окружения, синтезе объёмного контента~\cite{eigen2014depth,ranftl2020robust}. Современные архитектуры монокулярной оценки глубины обеспечивают высокое пространственное разрешение выходной карты и устойчивое восстановление относительной геометрии сцены. Вместе с тем выход таких моделей, как правило, определён с точностью до монотонного преобразования и масштаба (аффинная или более общая неопределённость относительно метрической глубины), что ограничивает прямое применение результатов в задачах, где требуются абсолютные расстояния в метрах или согласованный масштаб между кадрами.

Разрежённые измерения глубины --- типичный источник метрической информации: они поступают от ли\-даров, стереопар, структурированной подсветки или малоразмерных массивов прямого измерения времени пролёта (dToF). Если разрешение подсказки существенно ниже разрешения RGB, ассимилировать её с монокулярной оценкой можно как задачу \emph{условной калибровки}: сохранить пространственную структуру карты $d_{\mathrm{rel}}$ и согласовать её с разрежёнными наблюдениями. Альтернативная постановка --- полное depth completion --- предполагает обучение крупной сети на больших корпусах; в настоящей работе изучается более лёгкая семья методов, допускающая обучение ограниченного числа параметров на одном кадре или явную геометрическую параметризацию по областям изображения.

\textbf{Объект исследования} --- процедуры восстановления метрической карты глубины $\hat D$ по паре $(d_{\mathrm{rel}}, d_{\mathrm{sparse}}, M_{\mathrm{sparse}}, I)$ при фиксированной монокулярной модели.

\textbf{Предмет исследования} --- свойства и эффективность нескольких классов калибровочных преобразований (глобальное аффинное, локально-аффинное по суперпикселям с последующим сглаживанием, остаточная аппроксимация с помощью INR и условная INR-FiLM модель с постобработкой поля контекста).

\textbf{Новизна на уровне курсовой работы} заключается в систематическом сопоставлении указанных подходов на двух различных по происхождению приора протоколах (синтетическое разрежение NYU Depth~V2 и зональный приор набора ZJU-L5) при единой монокулярной основе Depth Anything~V2~\cite{yang2024depthanythingv2}.

\section{Цель и задачи работы}

\textbf{Цель работы} --- разработать и экспериментально сравнить методы восстановления метрической карты глубины по монокулярному предсказанию и разрежённому приору при ограниченной размерности обучаемой части на один кадр.

\textbf{Задачи работы:}
\begin{enumerate}
  \item формализовать постановку задачи и критерии качества в терминах согласованности $\hat D$ с эталоном $D_{\mathrm{gt}}$ на множестве валидных пикселей;
  \item описать и классифицировать реализованные методы: глобальная и локальная аффинная калибровка, INR на координатных признаках, INR-FiLM с региональным контекстом и регуляризующими процедурами на поле контекста;
  \item подобрать экспериментальные протоколы на наборах NYU Depth~V2 и ZJU-L5 с учётом различий в происхождении разрежённого приора;
  \item получить численные оценки качества и интерпретировать различия методов на агрегированных метриках.
\end{enumerate}

\section{Краткий обзор связанных работ}

Оценка глубины по одному изображению активно развивается как с учителем на помеченных корпусах~\cite{eigen2014depth}, так и в постановках устойчивости к переносу между доменами~\cite{ranftl2020robust}. Архитектура Depth Anything~V2~\cite{yang2024depthanythingv2} относится к семейству масштабированных трансформерных предикторов относительной глубины; её выход рассматривается в работе как фиксированный источник $d_{\mathrm{rel}}$, что позволяет изолировать вклад этапа калибровки.

Разбиение изображения на суперпиксели методом SLIC~\cite{achanta2012slic} широко применяется для локальной адаптации моделей к данным на областях со сходной цветовой и текстурной статистикой. Условная модуляция признаков по внешнему вектору (FiLM)~\cite{perez2018film} позволяет вводить низкоразмерное описание контекста без изменения общей архитектуры основной сети. Сохранение резких границ при сглаживании полей на изображении достигается, в частности, фильтрами, учитывающими направляющее изображение~\cite{he2010guided}.

\section{Постановка задачи}

Фиксируется один кадр. Заданы:
\begin{itemize}
  \item RGB-изображение $I\in\mathbb{R}^{H\times W\times 3}$;
  \item карта монокулярной относительной глубины $d_{\mathrm{rel}}\in\mathbb{R}^{H\times W}$, полученная предобученной моделью (положительность и предварительная нормализация определяются реализацией инференса);
  \item разрежённая метрическая подсказка $d_{\mathrm{sparse}}$ и маска наблюдений $M_{\mathrm{sparse}}\in\{0,1\}^{H\times W}$, согласованные с экспериментальным протоколом (синтетическое подвыборочное маскирование GT или зональные измерения датчика).
\end{itemize}

Пусть $\Omega_{\mathrm{gt}}=\{(x,y): D_{\mathrm{gt}}(x,y)>0\}$ --- множество пикселей, на которых определена эталонная метрическая глубина $D_{\mathrm{gt}}$ для оценки качества (маска достоверности задаётся данными и может совпадать или отличаться от носителя приора в зависимости от постановки).

\textbf{Задача}: построить отображение
\[
  \mathcal{F}:\ (d_{\mathrm{rel}}, d_{\mathrm{sparse}}, M_{\mathrm{sparse}}, I)\ \mapsto\ \hat D,
\]
минимизирующее в среднем по выборке кадров стандартные функционалы ошибки между $\hat D$ и $D_{\mathrm{gt}}$ на $\Omega_{\mathrm{gt}}$ (конкретные метрики приведены в разделе~\ref{sec:metrics}). Все рассматриваемые $\mathcal{F}$ параметризуются конечномерным вектором, оптимизация которого для INR осуществляется градиентным методом по разрежённому множеству $M_{\mathrm{sparse}}$.

\section{Методология}

\subsection{Глобальная аффинная калибровка}

Базовый класс преобразований
\[
  \hat D = s\, d_{\mathrm{rel}} + t,\qquad s>0,\quad t\in\mathbb{R}.
\]
Коэффициенты $(s,t)$ определяются методом наименьших квадратов (или эквивалентной $L_2$-минимизацией) по пикселям, где $M_{\mathrm{sparse}}=1$ и известны сопоставимые значения приора. Данный слой задаёт нижнюю границу сложности и единый масштаб по всему полю зрения.

\subsection{Локально аффинная калибровка по суперпикселям}

Изображение $I$ сегментируется на непересекающиеся области $\{\mathcal{R}_k\}_{k=1}^K$ методом SLIC~\cite{achanta2012slic}. На каждой области независимо ищется пара $(s_k,t_k)$ по тем же разрежённым наблюдениям; при недостатке точек используется глобальная оценка $(s,t)$ как априорное значение. Полученные кусочно-постоянные поля $(s(x,y),t(x,y))$ интерпретируются как низкочастотная аппроксимация пространственной зависимости масштаба и смещения.

Для подавления артефактов «ступенек» на границах $\mathcal{R}_k$ выполняется совместное билатеральное сглаживание пар $(s,t)$: пространственное ядро комбинируется с ядром по близости значений параметров, что ослабляет сглаживание через истинные разрывы глубины и сохраняет гладкость внутри однородных участков.

\subsection{Остаточная INR-модель на координатных признаках}

Вводится отображение пикселя $(x,y)$ в вектор признаков: нормированные координаты, синусно-косинусное разложение (positional encoding) по двум осям и нормированное значение $d_{\mathrm{rel}}(x,y)$. Небольшая многослойная перцептронная сеть $g_\theta$ аппроксимирует остаток $\Delta$ к глобальной аффинной базе:
\[
  \hat D(x,y) = s\, d_{\mathrm{rel}}(x,y) + t + \Delta_\theta(x,y).
\]
Параметры $\theta$ находятся минимизацией $\sum_{(x,y)\in M_{\mathrm{sparse}}} |\hat D(x,y)-d_{\mathrm{sparse}}(x,y)|$. Такая постановка интерпретируется как поиск гладкой поправки к базовой линии при ограниченном числе наблюдений. Плотное восстановление $\hat D$ на всей сетке достигается одним проходом сети после оптимизации.

\subsection{Условная INR-FiLM модель с региональным контекстом}

Каждая область SLIC сопоставляется вектору контекста $\mathbf{c}_k\in\mathbb{R}^{d_c}$, получаемому либо свёрточным кодированием локальной обрезки $(I, d_{\mathrm{rel}})$, либо комбинацией признаков замороженного ResNet18 и скалярного статистического признака по $d_{\mathrm{rel}}$ на области. Для пикселя $(x,y)\in\mathcal{R}_k$ входной вектор пиксельных признаков $\boldsymbol{\phi}(x,y)$ модулируется слоями FiLM~\cite{perez2018film}:
\[
  \mathbf{h}^{(\ell+1)} = \sigma\!\bigl(\mathbf{W}^{(\ell)}\mathbf{h}^{(\ell)} \odot (\mathbf{1}+\tanh \boldsymbol{\Gamma}^{(\ell)}(\mathbf{c}_k)) + \boldsymbol{\beta}^{(\ell)}(\mathbf{c}_k)\bigr),
\]
где $\boldsymbol{\Gamma}^{(\ell)}, \boldsymbol{\beta}^{(\ell)}$ --- линейные отображения контекста; $\sigma$ --- нелинейность; структура повторяется для $\ell=1,\ldots,L$. Выходной слой формирует остаток к глобальной аффинной базе аналогично п.~4.3.

Кусочно-постоянство $\mathbf{c}_k$ по областям приводит к разрывам модульных коэффициентов на границах SLIC. Для ослабления данного эффекта строится плотное поле $\mathbf{C}(x,y)$ размерности $(H,W,d_{\mathrm{tot}})$ путём назначения $\mathbf{c}_k$ всем пикселям области и последующего пространственного сглаживания: изотропная фильтрация Гаусса либо guided filter~\cite{he2010guided} с опорой на яркость $I$, что сохраняет контраст на истинных перепадах интенсивности. Идентичная процедура применяется при обучении и при выводе, что исключает рассогласование режимов. Дополнительно в $\mathbf{c}_k$ может конкатенироваться тройка геометрических дескрипторов области (нормированные координаты центроида, логарифм площади); в функцию потерь опционально вводится взвешенная по RGB-градиенту полная вариация на случайных рёбрах решётки для подавления шума при сохранении границ.

\subsection{Подход INR + softblend + refined}

\textbf{Идея.} Подход относится к классу лёгкой покадровой калибровки: вместо обучения крупной сетевой модели на датасете обучается компактная INR-параметризация для одного кадра, а пространственная согласованность обеспечивается мягким смешиванием локальных параметров (\textit{softblend}) и последующим уточнением (\textit{refined}) с сохранением границ.

\textbf{Формализация.} Пусть $\{(s_k,t_k)\}_{k=1}^K$ --- локальные аффинные параметры, оценённые на областях SLIC. Для пикселя $(x,y)$ итоговые параметры определяются как
\[
  s^\ast(x,y)=\sum_{k=1}^{K} w_k(x,y)\,s_k,\qquad
  t^\ast(x,y)=\sum_{k=1}^{K} w_k(x,y)\,t_k,\qquad
  \sum_{k=1}^{K} w_k(x,y)=1,\ \ w_k\ge 0,
\]
где веса $w_k$ вычисляются по близости к центрам областей и фотометрической согласованности (мягкое смешивание без жёстких границ сегментации). Далее INR-модель аппроксимирует остаток:
\[
  \hat D(x,y)=s^\ast(x,y)\,d_{\mathrm{rel}}(x,y)+t^\ast(x,y)+\Delta_\theta(x,y).
\]
Этап \textit{refined} использует edge-aware сглаживание (guided/bilateral) и/или слабую TV-регуляризацию для подавления локального шума в $\Delta_\theta$ при сохранении резких перепадов глубины.

\textbf{Режим оптимизации.} Параметры $\theta$ и (опционально) параметры смешивания подбираются по разрежённым наблюдениям $M_{\mathrm{sparse}}$; базовая глобальная аффинная оценка используется как стабилизирующая инициализация. Такой режим обеспечивает малую вычислительную стоимость и адаптацию под конкретный кадр.

\textbf{Ограничения.} Качество зависит от корректности разрежённого приора и устойчивости разбиения SLIC; при сильно неоднородных сценах возможно локальное переобучение остатка при недостатке измерений.

\subsection{Подход Direct-Depth CNN}

\textbf{Идея.} В отличие от покадровой калибровки, \textit{Direct-Depth CNN} рассматривает задачу как прямое отображение
\[
  \mathcal{G}_\psi:\ (I,\ d_{\mathrm{sparse}},\ M_{\mathrm{sparse}})\ \mapsto\ \hat D,
\]
где параметры $\psi$ общей сверточной сети обучаются на большом корпусе пар «вход--эталон».

\textbf{Архитектурный принцип.} Как правило, используется encoder--decoder с пропускными связями; RGB-ветвь извлекает семантику и структуру сцены, а ветвь разрежённой глубины передаёт метрический якорь. Слияние признаков выполняется на нескольких масштабах, после чего сеть восстанавливает плотную карту глубины.

\textbf{Сильные стороны.} Подход хорошо моделирует сложные нелинейные зависимости и при наличии репрезентативного обучения обеспечивает высокое среднее качество на целевом домене.

\textbf{Ограничения.} Требуются большие размеченные наборы и значительные вычислительные ресурсы обучения/инференса; перенос между доменами может сопровождаться деградацией. В рамках данной курсовой Direct-Depth CNN используется как концептуальный класс методов для сравнения постановок, тогда как практические эксперименты сосредоточены на лёгких калибровочных схемах поверх фиксированного монокулярного предсказания.

\section{Экспериментальная база и программная реализация}

\subsection{Набор NYU Depth V2}

Стандартный корпус пар RGB\,--\,метрическая глубина для внутренних сцен~\cite{eigen2014depth}. В экспериментах использовалось подмножество кадров из официальной разметки; разрежённый приор формировался случайным маскированием эталонной карты с заданной долей $\rho$ валидных пикселей (\texttt{sparse\_density}${}=\rho$) при фиксированном зерне генератора для воспроизводимости. Такой протокол моделирует идеализированный датчик с произвольным расположением точек измерения на высоком разрешении.

\subsection{Набор ZJU-L5}

Корпус реальных записей с RGB разрешения порядка $480\times640$ и низкоразрешённой подсказкой, соответствующей массиву VL53L5; эталонная глубина получена внешним высокоразрешённым датчиком в постановке набора. Использовалось тестовое разбиение; разрежённый приор извлекался из зональных данных датчика (\texttt{l5\_zones}), что отражает физически обусловленную структуру наблюдений и приближает задачу к конфигурации мобильных модулей глубины.

\subsection{Монокулярный блок}

Для всех экспериментов зафиксирована модель \texttt{depth-anything/Depth-Anything-V2-Small-hf}~\cite{yang2024depthanythingv2}; варианты калибровки отличаются только последующим отображением $\mathcal{F}$.

\subsection{Воспроизводимость}

Конфигурации экспериментов описываются декларативно (набор данных, число кадров, $\rho$, перечень методов и гиперпараметров); вычислительный эксперимент выполняется единым сценарием, результаты агрегируются в машиночитаемом виде и усредняются по кадрам выборки.

\section{Критерии качества}
\label{sec:metrics}

Ошибки вычисляются только на $\Omega_{\mathrm{gt}}$; ниже $p=\hat D$, $g=D_{\mathrm{gt}}$:
\begin{align*}
  \mathrm{AbsRel} &= \frac{1}{|\Omega_{\mathrm{gt}}|}\sum_{\Omega_{\mathrm{gt}}}\frac{|p-g|}{g}, &
  \mathrm{SqRel} &= \frac{1}{|\Omega_{\mathrm{gt}}|}\sum_{\Omega_{\mathrm{gt}}}\frac{(p-g)^2}{g},\\
  \mathrm{RMSE} &= \sqrt{\frac{1}{|\Omega_{\mathrm{gt}}|}\sum_{\Omega_{\mathrm{gt}}}(p-g)^2}, &
  \mathrm{RMSE}_{\log} &= \sqrt{\frac{1}{|\Omega_{\mathrm{gt}}|}\sum_{\Omega_{\mathrm{gt}}}\bigl(\ln\max(p,\varepsilon)-\ln\max(g,\varepsilon)\bigr)^2}.
\end{align*}
Пороговые метрики точности определяются как
\[
  \delta_k = \frac{1}{|\Omega_{\mathrm{gt}}|}\sum_{\Omega_{\mathrm{gt}}}\mathbf{1}\!\left[\max\!\left(\frac{p}{g},\frac{g}{p}\right)<1{,}25^k\right],\quad k\in\{1,2,3\}.
\]
Данный набор согласован с распространённой практикой оценки методов глубины~\cite{eigen2014depth}.

\section{Результаты и их интерпретация}

Приводятся агрегированные по всем кадрам выборки средние значения метрик (поле \texttt{summary} артефактов эксперимента); они характеризуют ожидаемое качество метода при случайном кадре из протокола. Абсолютная величина ошибки зависит от распределения глубин и масштаба сцены; сопоставимы между собой только методы внутри одной строки протокола.

\subsection{Протокол NYU Depth V2}

Выборка: $200$ кадров; $\rho=0{,}3$; сравнивались локально-аффинная калибровка с билатеральным сглаживанием $(s,t)$ и INR-FiLM при конфигурации прогона \texttt{inr\_film\_vs\_bilateral\_nyu200}. Результаты представлены в табл.~\ref{tab:nyu}.

\begin{table}[h]
  \centering
  \begin{tabular}{lcccccc}
    \toprule
    Метод & $\delta_1$ & $\delta_2$ & $\delta_3$ & AbsRel & RMSE \\
    \midrule
    INR-FiLM & 0{,}998 & 1{,}000 & 1{,}000 & 0{,}009 & 0{,}056 \\
    Локальная аффинная + bilateral & 0{,}986 & 0{,}997 & 0{,}999 & 0{,}024 & 0{,}157 \\
    \bottomrule
  \end{tabular}
  \caption{Средние метрики качества, NYU Depth V2.}
  \label{tab:nyu}
\end{table}

\textbf{Интерпретация.} При случайном разрежении высокого разрешения INR-FiLM достоверно снижает AbsRel и RMSE относительно локальной аффинной модели; высокие $\delta_k$ указывают на то, что большая доля пикселей попадает в узкий коридор относительных отклонений от GT. Это согласуется с гипотезой о необходимости нестационарной поправки к базовой линии, выходящей за пределы кусочно-постоянного $(s_k,t_k)$.

\subsection{Протокол ZJU-L5}

Выборка: $100$ кадров тестового множества; приор \texttt{l5\_zones}; дополнительно задавалось $\rho=0{,}3$ там, где протокол предполагает синтетическое разрежение совместно с зонами (конфигурация прогона \texttt{zju\_l5\_inr\_vs\_bilateral\_test}). Результаты --- табл.~\ref{tab:zju}.

\begin{table}[h]
  \centering
  \begin{tabular}{lcccccc}
    \toprule
    Метод & $\delta_1$ & $\delta_2$ & $\delta_3$ & AbsRel & RMSE \\
    \midrule
    INR-FiLM & 0{,}857 & 0{,}919 & 0{,}945 & 0{,}109 & 0{,}693 \\
    Локальная аффинная + bilateral & 0{,}857 & 0{,}920 & 0{,}947 & 0{,}108 & 0{,}704 \\
    \bottomrule
  \end{tabular}
  \caption{Средние метрики качества, ZJU-L5.}
  \label{tab:zju}
\end{table}

\textbf{Интерпретация.} Средние значения метрик для двух методов статистически близки на уровне округления; доминирующий вклад в ошибку вносит подмножество сложных кадров и возможное расхождение между приором и эталоном в постановке набора. Различия проявляются преимущественно локально (анализ карт ошибок и визуальное сравнение предсказаний). Дополнительные абляции по полю FiLM-контекста (гауссово сглаживание, guided filter, геометрические признаки, TV-регуляризация) описаны в конфигурации \texttt{zju\_l5\_inr\_film\_v2} и позволяют отделить вклад устранения артефактов SLIC от вклада самой FiLM-архитектуры.

\begin{figure}[h]
  \centering
  \fbox{\begin{minipage}{0.88\textwidth}\centering\vspace{1.8cm}
    \textit{Рисунки: примеры карт глубины и ошибок $|\hat D-D_{\mathrm{gt}}|$ по результатам экспериментов.}
  \vspace{1.8cm}\end{minipage}}
  \caption{Иллюстративный материал (формируется по данным каталога результатов эксперимента).}
  \label{fig:placeholder}
\end{figure}

\section{Заключение}

В курсовой работе формализована задача восстановления метрической карты глубины при фиксированном монокулярном предсказании и разрежённом приоре; описана иерархия методов от глобальной аффинной подгонки к локально-аффинным моделям по областям SLIC и к остаточным INR-представлениям с условной FiLM-модуляцией по региональному контексту. Показано, что пространственное сглаживание параметров и полей контекста необходимо для подавления артефактов сегментации без размывания истинных границ при использовании guided фильтрации~\cite{he2010guided}.

Экспериментально установлено существенное преимущество INR-FiLM перед локальной аффинной моделью на протоколе NYU Depth~V2 при случайном разрежении фиксированной плотности; на протоколе ZJU-L5 с зональным приором средние метрики оказались сопоставимыми, что указывает на необходимость более тонкого анализа по подмножествам кадров и уточнения постановки при использовании реальных низкоразрешённых измерений.

\textbf{Направления дальнейшего развития}: статистический анализ значимости различий по кадрам; расширение семейства сравниваемых методов; ужесточение связи с постановками depth completion на общих бенчмарках; исследование устойчивости к шуму и ошибке калибровки приора.

\textbf{Практическая значимость} результатов заключается в возможности встраивания лёгких процедур калибровки поверх готовых монокулярных моделей в приложениях с ограниченными вычислительными ресурсами и частичными измерениями глубины.

\bibliography{refs_coursework}

\end{document}
